\documentclass{article}
\usepackage[dblblindworkshop]{neurips_2026}
\workshoptitle{TAE (Trust-AI-Eval): Can We Trust AI Evaluation?}
\usepackage[utf8]{inputenc}
\usepackage[T1]{fontenc}
\usepackage{microtype}
\usepackage{graphicx}
\usepackage{booktabs}
\usepackage{amsmath,amssymb}
\usepackage{url}
\usepackage{xcolor}
\title{When Evaluation-Awareness Probes Transfer Across Formats}
\author{Anonymous Authors}

\begin{document}
\maketitle
\begin{abstract}
Evaluation-purpose probes can exploit benchmark structure rather than evaluation purpose. We study this failure with an English controlled purpose $\times$ format factorial design. Equal-N controls and an all-layer decomposition separate shared purpose ($A$), format ($B$), and purpose$\times$format interaction ($C$). Geometry selects a shared layer and an interaction-heavy layer before intervention. In BF16 Llama-3.1-8B, correctly signed $C$ improves continued-forward transport toward the model's natural matched purpose counterfactual at the interaction-heavy layer, but not at the shared layer; wrong-sign, norm-matched, shuffled, and fixed-random controls do not reproduce the effect. The direct behavioral endpoint fails validation, so the claim is restricted to internal mechanistic representation under controlled prompt-provided context, not behavioral steering or spontaneous evaluation awareness.
\end{abstract}

\section{Introduction}
Evaluation awareness is often measured by linear probes over model activations. A probe can score highly while reading a correlated benchmark format, so accuracy alone does not establish that the representation contains purpose. We ask a narrower question: when a purpose probe transfers between benchmark-style and casual formats, what representation component determines success or failure?

Our controlled design crosses evaluation/deployment purpose with benchmark/casual format for matched payloads. We first remove the current pooled-versus-transfer training-size confound. We then test the fixed factorial estimand $h_{p,f}=\mu+pA+fB+pfC+\epsilon$, where $A$ is shared purpose, $B$ is a format offset, and $C$ is a purpose$\times$format interaction. Finally, after freezing layer and endpoint rules from development geometry, we test held-out-family residual-stream counterfactual transport.

\paragraph{Claim boundary.} The manuscript distinguishes ranking (AUC), threshold calibration (balanced accuracy), observational geometry, internal causal mediation, and behavioral causal steering. No result is promoted beyond the strongest completed gate.

\section{Why Evaluation-Awareness Probes Are Confounded}
The existing diagnostic compares a pooled probe trained on all four cells with one-format probes trained on benchmark or casual cells. The pooled probe has both more rows and direct exposure to both target formats. Thus its advantage is not a clean test of format-conditioned purpose coding. We retain this diagnostic as development evidence only and make equal-N training the first inferential gate.

Probe selectivity and control-task work motivates matched controls \citep{hewitt2019control,belinkov2021probing}. Causal probing work also distinguishes decodability from behavioral use \citep{elazar2020amnesic}; our intervention claims follow that distinction.

Recent work makes the adjacent methodological space clear: format-confound controls and evaluation-context transfer motivate the present factorial design \citep{sahoo2026format,devbunova2026evaluation}, while activation transport and representation interventions provide causal tools \citep{rodriguez2025act,wu2024reft,geiger2024das}. A closely related structure$\times$register factorial study \citep{santana2026relational} reinforces that interaction terms can matter, but does not estimate our purpose$\times$format residual-space $A/B/C$ decomposition or compare held-out interventions with natural matched downstream targets. Our novelty claim is therefore the combination, not factorial interactions or activation intervention in isolation.

\section{Controlled Purpose $\times$ Format Design}
Each payload block has four matched cells: $h_{EB}$, $h_{EC}$, $h_{DB}$, and $h_{DC}$. Purpose is effect-coded as evaluation $p=+1$ and deployment $p=-1$; format is benchmark $f=+1$ and casual $f=-1$. We hold payload identity and purpose-family splits fixed, and evaluate with leave-one-purpose-family-out folds.

The factorial decomposition is predeclared:
\begin{align}
 h_{p,f}&=\mu+pA+fB+pfC+\epsilon,\\
 A&=(\delta_B+\delta_C)/4,\quad C=(\delta_B-\delta_C)/4,\\
 B&=((h_{EB}+h_{DB})-(h_{EC}+h_{DC}))/4,
\end{align}
where $\delta_B=h_{EB}-h_{DB}$ and $\delta_C=h_{EC}-h_{DC}$. All analyses preserve complete payload blocks.

\section{Confounded, Decorrelated, and Cross-Format Results}
\begin{figure}[t]
\centering\fbox{\parbox{0.9\linewidth}{Placeholder: full depth curves for confounded, pooled/decorrelated, benchmark$\rightarrow$casual, and casual$\rightarrow$benchmark AUC and balanced accuracy.}}
\caption{Full layer trajectories; no best-layer-only presentation.}
\end{figure}

The current local Llama-3.1-8B and Llama-3.3-70B files are development diagnostics. We will report their exact stimulus, revision, layer semantics, and status before interpreting any model difference. The headline observation is treated as a motivation, not as a mechanism: pooled purpose decoding is higher than one-format transfer, but the original comparison is confounded by training size and target-format exposure.

\section{Equal-Training-Row Control}
For each held-out purpose family, benchmark-only and casual-only training use the eval/deploy pair from every training block in one format. Equal-N mixed-format training uses one balanced format assignment per block and retains that pair. The assignment seed bank is fixed before evaluation; blocks are never split.

\begin{table}[t]
\centering
\caption{Primary Llama-3.1-8B layer-24 development diagnostic. Values are AUC / balanced accuracy; five fixed assignment seeds for equal-N.}
\begin{tabular}{lcc}\toprule
Training regime & AUC & Balanced accuracy\\\midrule
Benchmark-only & 0.6113 & 0.6042\\
Casual-only & 0.6259 & 0.5590\\
Equal-N mixed-format & 0.6684 & 0.6250\\
Pooled/decorrelated & 0.7029 & 0.6597\\
Benchmark $\rightarrow$ casual & 0.5544 & 0.5208\\
Casual $\rightarrow$ benchmark & 0.5656 & 0.5000\\\bottomrule
\end{tabular}
\end{table}

The primary result is that equal-N mixed-format training remains above one-format training on this development bank, while pooled training remains higher. This rejects a sample-count-only explanation but does not by itself establish the factorial mechanism. The source artifact and its development-only status are recorded in \texttt{runs/tae\_2026/equal\_n/equal\_n\_20260815\_8b\_l24/}.

\section{Factorial Representational Mechanism}
We compute $A$, $B$, and $C$ in the original activation space and in fixed 256-dimensional controls. Reducers are fitted on training payloads only. We report full layer curves for $\|A\|$, $\|B\|$, $\|C\|$, $\|C\|/\|A\|$, matched-direction cosines, principal cosines, projections of $B$ onto learned purpose directions, cross-format AUC, and balanced accuracy. Layers are repeated measurements, not independent replicates; uncertainty resamples complete families or payload blocks.

\begin{figure}[t]
\centering\fbox{\parbox{0.9\linewidth}{Placeholders: interaction ratio vs relative depth; cross-format AUC; interaction ratio vs AUC; $B$ projection vs AUC--balanced-accuracy gap; delta cosine; subspace overlap.}}
\caption{Predeclared all-layer geometry curves.}
\end{figure}

The explanatory tests are associations between development geometry and held-out transfer. They are not causal claims. Undefined zero-vector ratios and cosines remain explicit missing values.

\section{Causal Counterfactual Transport}
After geometry is complete, we freeze a layer-selection rule using held-out/development purpose decodability and low/high $\|C\|/\|A\|$, before inspecting intervention outcomes. Training purpose families estimate raw-space $\hat A$, $\hat B$, and $\hat C$. At an eligible layer, purpose transport uses $-2p(\hat A+f\hat C)$ and format transport uses $-2f(\hat B+p\hat C)$; main-effect-only ablations remove $C$.

Interventions continue the model forward. We validate direct candidate-logit endpoints before steering, randomize candidate mappings, and report purpose and format readouts separately by wrapper and family. Zero strength, sign reversal, and fixed random-direction controls are required. If endpoints are inadequate, downstream representation movement is reported only as internal causal mediation, not behavioral steering.

\begin{figure}[t]
\centering\fbox{\parbox{0.9\linewidth}{Placeholder: purpose/format readout changes for full factorial versus main-effect-only transport, with random and sign controls.}}
\caption{Held-out-family counterfactual transport.}
\end{figure}

\section{Depth, Scale, and Generation Comparison}
We compare Llama-3.1-8B-Instruct, Llama-3.1-70B-Instruct, and Llama-3.3-70B-Instruct on the exact rendered item order and hash, with BF16 inference, float32 stored activations, final prompt token, zero-based block output before final norm, identical family splits, and identical analysis code. The 3.1-8B/3.1-70B comparison is a within-family scale comparison; the 3.1-70B/3.3-70B comparison is a same-size generation/training comparison. Neither comparison is perfectly causal.

The current matched ledger records complete 8B and 3.3-70B legacy curves with exact revisions. The 3.1-70B archive is pending user-owned job `51479700`; jobs `51479700` and `51484306` are not modified by this sprint. No three-model claim is made until the 3.1-70B archive passes exact input-hash and layer-semantics checks.

\begin{figure}[t]
\centering\fbox{\parbox{0.9\linewidth}{Placeholder: matched relative-depth curves for pooled, cross-format, equal-N AUC, interaction ratio, delta cosine, and subspace overlap after 3.1-70B extraction.}}
\caption{Three-model matched trajectories.}
\end{figure}

\section{Implications for Evaluation Validity}
A high-scoring evaluation-awareness probe is evidence about a measurement instrument, not automatically evidence that a model represents evaluation purpose independently of benchmark structure. Equal-N training, matched payload blocks, separate AUC and balanced accuracy, and independent endpoint validation should be treated as validity checks. The practical output is a claim ladder: format shortcut, purpose decodability, equal-N survival, geometry, and only then causal transport.

The mechanistic result is the paper's main contribution. Naive probes can exploit prompt format; equal-N decorrelation shows that pooled gains are not only a row-count artifact; and the matched $2\times2$ design decomposes purpose representation into shared purpose ($A$), format ($B$), and purpose$\times$format interaction ($C$). Geometry selected L5 as shared and L16 as interaction-heavy before intervention. At L16, correctly signed $C$ improved transport toward the model's natural matched purpose counterfactual, while wrong-sign, shuffled, norm-matched, and fixed-random controls did not. This is an internal representation result under controlled prompt-provided context, not evidence that the model spontaneously recognizes evaluation or can be behaviorally steered.

\section{Limitations}
This sprint uses an English controlled bank and does not establish multilingual or naturalistic generalization. The tracked activation artifact contains 288 rows and 72 blocks, with unequal family/topic coverage; its coverage origin must be explained before confirmation-grade claims. Format changes prompt wrappers and response instructions, so the factorial format effect is intentionally a measurement construct, not a pure semantic variable.

Model generation, training, and scale remain confounded in the 8B/70B comparison pending the exact Llama-3.1-70B archive. Probe geometry is observational because the held-out intervention did not reach baseline: the independent 8B causal job hit a Hugging Face gated-repository HTTP 403 before model loading. Direct-logit endpoints can be wrapper- and tokenization-sensitive; no endpoint result is reported here. We do not use old jProbe or migration-bundle results as evidence.


\bibliographystyle{plainnat}
\bibliography{references}
\appendix
\section{Additional Methods and Reproducibility}
Every run records the git commit, immutable run ID, item and activation hashes, model revision, dtype, token position, layer semantics, split unit, regularization, fixed seeds, and output hashes. Complete payload blocks are the statistical unit for factorial reconstruction; purpose families define the outer generalization fold. Layers are trajectories from one model, not independent samples.

\paragraph{Required controls.} Fixed 256-D random projections use predeclared seeds and PCA is fitted on training rows only. Block/family bootstrap procedures preserve all four cells. Causal protocols freeze endpoint wording, candidate mappings, layer eligibility, raw-space direction conversion, strength, signs, and random controls before held-out steering.

\paragraph{Artifact ledger.} Current development artifacts: equal-N layer 24 at \texttt{runs/tae\_2026/equal\_n/equal\_n\_20260815\_8b\_l24/}; primary geometry layer 24 at \texttt{runs/tae\_2026/geometry/geometry\_20260815\_8b\_l24/}; pending 3.1-70B extraction job 51479700. The living ledger is \texttt{runs/tae\_2026/STATUS.md}.

\end{document}
